\section{Introduction}

Forest harvesting operations must coordinate machines, crews, and harvest blocks over days and
weeks while staying consistent with tactical and strategic plans \citep{martell1998forest}.
Delays in deploying equipment disrupt wood supply to mills and raise operating costs
\citep{d2008using,alayet2013planning}, particularly where weather, breakdowns, and markets make
operating conditions uncertain \citep{arora2023detailed}. Machine scheduling, the assignment of
machines to blocks over discrete time periods subject to spatial, temporal, and logistical
constraints \citep{corner2005scheduling}, is therefore a core operational planning problem.
Practice ranges from spreadsheets and whiteboards in small and medium operations to proprietary
optimization software in large ones. Few of these tools are open to inspection, and design choices
in planning systems are rarely evaluated systematically. As a result, findings about how to
configure them are hard to reproduce or transfer \citep{jaffray2026thesis}.

Machine scheduling problems are large and combinatorial \citep{bellavenutte2020partitioning}.
Deterministic optimization over long horizons is often intractable, and long deterministic
schedules lose value as conditions change \citep{sethi1991theory}. Rolling-horizon planning
addresses both problems. It solves a sequence of shorter subproblems, implements (locks) the first
part of each solution, advances time, and re-optimizes \citep{baker1979analytic,
venkataraman1999effect}. In forest operations this lets near-term decisions be fixed while later
decisions stay flexible \citep{feng2010simulation,beaudoin2008hierarchical}.

Rolling-horizon performance depends on two design parameters: the length of each planning window
(the planning horizon) and the interval between re-optimizations, which is also the span of
decisions committed at each step. Operations research theory suggests a horizon beyond which
further look-ahead adds little \citep{baker1979analytic,sethi1991theory,santos2021integrated}.
In practice, horizons are often set by convention, which can impose high computing costs or rely
on unreliable long-range forecasts \citep{chand2002forecast}. In forest machine scheduling there is
little evidence on how horizon length and re-optimization frequency trade off solution quality
against computing effort. Uncertainty is usually the motivation for re-planning, but the
structural effect of these parameters can be isolated under deterministic conditions first.

A sound evaluation also depends on how the rolling procedure is implemented and how its output is
scored. Each window must start from the state left by the decisions already implemented. In a
multi-phase harvesting system (felling, skidding, processing, loading), that state includes the
volume remaining in each block, the inventory staged between phases, and the position of each
machine. A rolling plan should be judged by the schedule that is actually implemented, that is,
the stitched sequence of locked decisions, evaluated against the full problem. The objective of
any single window is not a measure of that schedule. An implementation that omits the carried
state, or a metric based on a single window, can produce apparent horizon and frequency effects
that are artifacts. \citet{jaffray2026thesis} reported a first evaluation of these parameters with
an early release of the Forest Harvesting Operations Planning System (FHOPS). That release did not
carry state between windows, and the evaluation scored runs by the last window's objective. FHOPS
1.0.1 \citep{fhops2026v101} carries the operational state forward between windows and replays the
stitched plan for evaluation.

This study repeats that evaluation with FHOPS 1.0.1, an open-source machine scheduling framework
(\fhopspaper), on nine semi-synthetic case studies representative of British Columbia (BC)
harvesting operations. We ask:
\begin{enumerate}
  \item At what planning horizon length do further gains in delivered volume become small?
  \item Does frequent re-optimization outperform longer planning windows that are updated less
        often?
  \item How do horizon length and re-optimization frequency affect schedule continuity and the
        feasibility of the implemented schedule?
  \item How do these effects depend on problem size and on the solution method (exact
        mixed-integer programming versus simulated annealing), including the computing time
        each configuration requires?
\end{enumerate}
All experiments are deterministic. They therefore establish baseline performance bounds that
future studies under uncertainty can be compared with.
